\section{Họ xóa dần}
\label{sec:ch08-ho-xoa-dan}

\index{deletion}%
\index{insertion}%
\index{comprehensiveness}%
\index{sufficiency}%
Họ đầu tiên điền vào ô thứ nhất bằng cách tự nhiên nhất: không chọn một tập con
cố định, mà chạy hết mọi tập con theo thứ tự điểm và vẽ ra một đường cong.
Petsiuk và cộng sự đặt tên cho hai đường cong ấy khi giới thiệu phương pháp
RISE~\autocite{mrise}.

Đường cong thứ nhất, deletion, đo xác suất của lớp được dự đoán giảm đi thế nào
khi các điểm ảnh quan trọng lần lượt bị bỏ; các tác giả viết rằng một cú giảm
dốc, tức diện tích dưới đường cong nhỏ, nghĩa là lời giải thích tốt. Đường cong
thứ hai, insertion, làm ngược lại: nó đo xác suất tăng lên thế nào khi các điểm
ảnh lần lượt được đưa vào, và ở đây diện tích lớn mới là tốt~\autocite{mrise}.
Trực giác đằng sau deletion được nêu thẳng trong bài: bỏ đi \enquote{nguyên
nhân} thì sẽ buộc mô hình gốc đổi quyết định~\autocite{mrise}.

Cùng ý tưởng ấy xuất hiện lại ở xử lý ngôn ngữ dưới hai cái tên khác, trong bộ
đánh giá ERASER của DeYoung và cộng sự. Ở đó lời giải thích là một tập đoạn văn
bản được đánh dấu, và hai chỉ số so độ tự tin của mô hình trước và sau khi động
vào tập ấy. Comprehensiveness lấy hiệu giữa độ tự tin trên đầu vào đầy đủ và độ
tự tin sau khi bỏ đi phần được đánh dấu, theo lập luận rằng mô hình phải bớt tự
tin khi phần ấy biến mất. Sufficiency lấy hiệu giữa độ tự tin trên đầu vào đầy
đủ và độ tự tin khi chỉ giữ lại phần được đánh dấu, tức đo phần ấy một mình đã
đủ cho một dự đoán hay chưa~\autocite{meraser}.

Bốn cái tên ấy đến từ hai bài báo ở hai lĩnh vực, nhưng chúng chạy cùng một
vòng lặp. So với
hình~\ref{fig:ch08-vong-lap-do}, deletion và comprehensiveness cùng chọn tập
đặc trưng điểm cao rồi bỏ đi, còn insertion và sufficiency cùng chọn tập điểm
cao rồi giữ lại một mình nó. Hai cách ấy đúng là hai đại lượng mà
chương~\ref{ch:khung-hop-nhat} đã nhận từ bài báo quan điểm, lần này có tên
riêng và có bài báo gốc đứng sau: giữ lại phần điểm cao là đo tính đủ, tức
fidelity, còn bỏ phần điểm cao đi là đo tính cần, tức inverse fidelity. Vậy
insertion và sufficiency thuộc về fidelity, deletion và comprehensiveness thuộc
về inverse fidelity.

Một đường cong deletion đủ ngắn thì tính được bằng tay, và ví dụ mà Phần~II đã
dùng cho đúng việc ấy. Vẫn $f(x_1,x_2) = x_1 + 2x_2 + x_1x_2$ tại $x = (1,1)$,
xóa nghĩa là đặt đặc trưng về 0. Các \tn{giá trị Shapley}{Shapley value} mà
mục~\ref{sec:ch04-tro-choi-chia-gia-tri} tính được xếp đặc trưng 2 trên đặc
trưng 1, nên thứ tự xóa là đặc trưng 2 rồi đặc trưng 1. Đầu ra lần lượt là
\[
  f(1,1) = 4, \qquad f(1,0) = 1, \qquad f(0,0) = 0 .
\]
Đặt trục ngang là phần đặc trưng đã xóa, tức $0$, $1/2$, $1$, rồi lấy diện tích
theo quy tắc hình thang:
\[
  \frac{4+1}{2}\cdot\frac{1}{2} \;+\; \frac{1+0}{2}\cdot\frac{1}{2}
  \;=\; 1.5 .
\]
Xóa theo thứ tự ngược lại cho đầu ra $4$, $2$, $0$ và diện tích $2.0$. Đây là
diện tích thô, chưa chia cho đầu ra ban đầu.

Hai con số ấy nói đúng một điều: thứ tự do Shapley đưa ra cho diện tích nhỏ
hơn thứ tự ngược, nên theo chiều đọc của deletion thì nó là thứ tự tốt hơn. Cái
chúng không nói là $1.5$ có phải một giá trị tốt hay không, vì không có thang
nào để so ngoài chính thứ tự ngược. Với hai đặc trưng thì chỉ có hai thứ tự và
phép so là đủ; với $d$ đặc trưng thì có $d!$ thứ tự và con số này chỉ so được
với những thứ tự mà ta chịu khó tính ra.

Ô thứ hai, giá trị thay thế, là chỗ họ này để lộ vấn đề của nó, và điều đáng
chú ý là chính các tác giả đã gặp vấn đề ấy nhưng không gọi nó là một lựa
chọn. RISE
không bắt đầu insertion từ một ảnh trống mà từ một ảnh đã làm mờ, với lý do được
nêu rõ: làm mờ lấy đi phần lớn chi tiết nhỏ mà không tạo ra các cạnh
sắc~\autocite{mrise}. Nói cách khác, các tác giả biết rằng một giá trị thay thế
vụng về sẽ tự nó sinh ra tín hiệu, nên họ chọn một giá trị ít sinh tín hiệu hơn.
Việc chọn ấy không được biện minh bằng gì ngoài chính nhận xét đó.

ERASER nêu giới hạn của mình ở một chỗ khác. Các tác giả viết rằng đo độ trung
thực của một đoạn được đánh dấu thế nào cho đúng vẫn là một câu hỏi mở, và rằng
trong phiên bản đầu tiên của bộ đánh giá này họ đề xuất những chỉ số đơn giản,
lấy cảm hứng từ công trình trước~\autocite{meraser}. Câu ấy đặt hai chỉ số vào
đúng vị trí của chúng: một điểm khởi đầu hợp lý, chứ không phải một dụng cụ đã
kiểm định.

Còn một điều mà cả bốn chỉ số đều không kiểm tra, và nó không nằm ở ô nào trong
ba ô. Đầu vào đã bị bỏ bớt đặc trưng là một đầu vào mà mô hình chưa từng thấy
khi huấn luyện. Mục~\ref{sec:ch08-huan-luyen-lai} đọc một bài báo đặt đúng vấn
đề đó và đề xuất một cách sửa.
